% Use the existing owner-derived source mechanism; do not edit frozen originals.
\OLSADeclareDocumentTextCorrection{OLP-0027}{FA0027-FINITE-ALT}
{تناظرهای دوسویی با $\Nat$ تعریف می‌کند.}
{تناظرهای دوسویی با $\Nat$ یا با بخشی ابتدایی از آن تعریف می‌کند.}
\OLSADeclareDocumentTextCorrection{OLP-0027}{FA0027-BOOK-TITLE}
{باتن برای استفاده در متن خود با عنوان «نظریهٔ مجموعه‌های باز» نوشته}
{باتن برای استفاده در کتاب خود با عنوان «Open Set Theory» نوشته}
\OLSADeclareDocumentTextCorrection{OLP-0027}{FA0027-ENUM-CONSISTENCY}
{شمارش‌پذیری را به کار می‌گیرند}
{شمارایی را به کار می‌گیرند}
\OLSADeclareDocumentTextCorrection{OLP-0027}{FA0027-INITIAL-SEGMENT}
{یک پارهٔ آغازین}
{یک بخش ابتدایی}
\OLSADeclareDocumentTextCorrection{OLP-0028}{FA0028-INFINITY-NOUN}
{یک نامتناهیِ
بالفعل}
{یک بی‌نهایتِ
بالفعل}
\OLSADeclareDocumentTextCorrection{OLP-0028}{FA0028-ENUM-GLOSS}
{مفهوم برشماری است.}
{مفهوم برشماری، یعنی فهرست‌کردن اعضا، است.}
\OLSADeclareDocumentTextCorrection{OLP-0029}{FA0029-INITIAL-SEGMENT}
{یک قطعهٔ آغازین}
{یک بخش ابتدایی}
\OLSADeclareDocumentTextCorrection{OLP-0029}{FA0029-FINITE-POSITION}
{  !!{element} داشته باشد که بی‌درنگ پیش از آن بیاید. به بیان دیگر، میان
  نخستین !!{element} آن فهرست و هر !!{element} دیگر فقط شمار متناهی‌ای عضو
  قرار دارد.}
{  !!{element} داشته باشد که بی‌درنگ پیش از آن بیاید. همچنین لازم است میان
  نخستین !!{element} آن فهرست و هر !!{element} دیگر فقط شمار متناهی‌ای عضو
  قرار داشته باشد. \emph{توضیح ویراستاری:} شرط اخیر از داشتن آغاز و عضو پیشینِ
  بلافصل به‌تنهایی نتیجه نمی‌شود؛ همان شرط جایگاه متناهی در تعریف بالاست.}
\OLSADeclareDocumentTextCorrection{OLP-0029}{FA0029-LIST-VS-FORMAL}
{اگر $A$ برشماری‌ای داشته باشد، برشماری‌ای بی‌تکرار نیز دارد.}
{اگر $A$ برشماری‌ای داشته باشد، برشماری‌ای بی‌تکرار نیز دارد.
  در این گزاره، «برشماری» به معنای فهرست در تعریف غیرصوری بالاست.}
\OLSADeclareDocumentTextCorrection{OLP-0029}{FA0029-NATURAL-EQUIVALENCE}
{بیایید خود را متقاعد کنیم که تعریف صوری و تعریف غیرصوری با استفاده از
فهرستی احتمالاً نامتناهی، هم‌ارزند.}
{برای مجموعهٔ ناتهی، هم‌ارزی تعریف صوری با تعریف غیرصوری بر پایهٔ
فهرستی احتمالاً نامتناهی را در هر دو جهت بررسی می‌کنیم.}
\OLSADeclareDocumentTextCorrection{OLP-0029}{FA0029-COUNTABLE-CONVENTION}
{مجموعهٔ~$A$ !!{enumerable} است اگر و تنها اگر تهی باشد یا برشماری‌ای داشته باشد.}
{مجموعهٔ~$A$ !!{enumerable} است اگر و تنها اگر تهی باشد یا برشماری‌ای داشته باشد.

\emph{یادداشت اصطلاح‌شناختی:} در این متن «شمارا» مجموعه‌های متناهی، از جمله
مجموعهٔ تهی، و مجموعه‌های شمارای نامتناهی را در بر می‌گیرد. برخی منابع فارسی
این واژه را فقط برای حالت نامتناهی به کار می‌برند؛ در اینجا تعریف متن مبدأ ملاک است.}
\OLSADeclareDocumentTextCorrection{OLP-0029}{FA0029-CEILING-EXACT}
{تابع \emph{سقف} را نشان می‌دهد که $x$ را رو به بالا
تا نزدیک‌ترین عدد صحیح گرد می‌کند}
{تابع \emph{سقف} را نشان می‌دهد که کوچک‌ترین عدد صحیحِ بزرگ‌تر یا مساویِ $x$ را
به دست می‌دهد}
\OLSADeclareDocumentTextCorrection{OLP-0029}{FA0029-MISSING-MINUS-THREE}
{0 & 1 & -1 & 2 & -2 & 3 & \dots}
{0 & 1 & -1 & 2 & -2 & 3 & -3 & \dots}
\OLSADeclareDocumentTextCorrection{OLP-0029}{FA0029-TABLE-DISCLOSURE}
{همچنین می‌توانید $f$ را به‌صورت موردیِ زیر در نظر بگیرید:}
{\emph{یادداشت ویراستاری:} در سطر آخر جدولِ متن مبدأ، مقدار هفتم، یعنی منفی سه،
جا افتاده بود؛ این مقدار در اینجا افزوده شده است.
همچنین می‌توانید $f$ را به‌صورت موردیِ زیر در نظر بگیرید:}
\OLSADeclareDocumentTextCorrection{OLP-0029}{FA0029-UNION-EMPTY-GRAMMAR}
{حالت‌هایی را که $A$ یا~$B = \emptyset$ است نیز بررسی کنید.}
{حالت تهی‌بودن $A$ و نیز حالت $B = \emptyset$ را بررسی کنید.}
\OLSADeclareDocumentTextCorrection{OLP-0029}{FA0029-RANGE-EXISTS-MEMBER}
{باید
  !!a{element} از~$A$ در برشماریِ}
{باید
  یک !!{element} از~$A$ در برشماریِ}
\OLSADeclareDocumentTextCorrection{OLP-0029}{FA0029-INDEFINITE-MEMBER}
{اگر !!a{element} از~$A$}
{اگر یک !!{element} از~$A$}
\OLSADeclareDocumentTextCorrection{OLP-0029}{FA0029-INFINITE-DOMAIN-GRAMMAR}
{بر بی‌نهایت !!{element}sی
مجموعهٔ~$\PosInt$}
{بر همهٔ !!{element}sی
مجموعهٔ نامتناهیِ~$\PosInt$}
\OLSADeclareDocumentTextCorrection{OLP-0029}{FA0029-POSITIVE-SQUARES}
{پنج عدد مربعی نخست}
{پنج مربعِ مثبتِ نخست}
\OLSADeclareDocumentTextCorrection{OLP-0030}{FA0030-PROOF-COMPLETENESS}
{مقدار خانهٔ واقع در سطر $n$اُم و ستون $m$اُمِ آرایهٔ زیگزاگ کانتور باشد.}
{مقدار خانهٔ واقع در سطر $n$اُم و ستون $m$اُمِ آرایهٔ زیگزاگ کانتور باشد.
\emph{تکمیل توضیح برهان:} خانه‌ها بر روی قطرهایی با مجموعِ ثابتِ دو مؤلفه
قرار دارند. قطرها به ترتیبِ این مجموع و خانه‌های هر قطر به ترتیبِ مؤلفهٔ اول
پیموده می‌شوند. هر قطر متناهی است و پیش از هر خانه فقط تعداد متناهی‌ای خانه
قرار می‌گیرد. این قطرها همهٔ زوج‌ها را دقیقاً یک‌بار در بر می‌گیرند؛ پس تابع
تعریف‌شده خوش‌تعریف و پوشاست.}
\OLSADeclareDocumentTextCorrection{OLP-0030}{FA0030-ZERO-POWER}
{بنابراین داریم:}
{\emph{یادداشت ویراستاری:} حالت توان صفر نیز شامل گزارهٔ بعدی است:
طبق قرارداد معمول، مجموعهٔ صفر‌تایی‌ها تک‌عضوی است و تنها عضو آن دنبالهٔ تهی
است؛ بنابراین این حالت نیز شماراست. شناسایی سه‌تایی با زوجِ شامل یک زوج،
همان قرارداد گروه‌بندی مؤلفه‌ها در متن مبدأ است.
بنابراین داریم:}
\OLSADeclareSetup{OLP-0027}{\olsa_source_correction_install_document_text:}
\OLSADeclareSetup{OLP-0028}{\olsa_source_correction_install_document_text:}
\OLSADeclareSetup{OLP-0029}{\olsa_source_correction_install_document_text:}
\OLSADeclareSetup{OLP-0030}{\olsa_source_correction_install_document_text:}
